% GENERATED FILE. Edit canonical lessons, then run npm run generate:books.
% canonical-source-hash: fnv1a64:b2c4cd5a0035ee54
% canonical-lessons: TA-C558-natasvaram, TA-C558-melam, TA-C558-rakam, TA-C558-talam, TA-C558-piyano, TA-R558-review-552-558, TA-R558-review-final-558

\chapter{Booking Musicians}
\label{ch:ta-booking-musicians}
\hlchaptermodality{558}

\begin{chapteropening}
\textbf{By the end of this chapter:} \emph{I can talk about booking musicians for a function.}

Complete the last lesson of chapter 558: I can talk about booking musicians for a function.
\end{chapteropening}

\section[\texorpdfstring{\ta{நாதஸ்வரம்} (nātasvaram)}{நாதஸ்வரம் (nātasvaram)}]{\ta{நாதஸ்வரம்} (nātasvaram) — a nagaswaram}

\label{lesson:TA-C558-natasvaram}



\begin{quote}
Before the new one: say the Tamil for a guitar, then the Tamil for a violin.
\end{quote}

\subsection*{You'll want to know: \ta{நாதஸ்வரம்}}
\textbf{\ta{நாதஸ்வரம்}} — \emph{nātasvaram} — \textquotedblleft{}a nagaswaram\textquotedblright{}.

\begin{quote}
\ta{நாதஸ்வரம்} \ta{இல்லாமல்} \ta{திருமணம்} \ta{நடக்காது} \ta{என்று} \ta{பாட்டி} \ta{சொல்கிறாள்}.
\end{quote}

\begin{quote}
\emph{nātasvaram illāmal tirumaṇam naṭakkātu eṉṟu pāṭṭi colkiṟāḷ.}
\end{quote}

\begin{quote}
Grandmother says a wedding cannot happen without the nagaswaram.
\end{quote}

A long double-reed wind instrument played at weddings and in temples.

\begin{cousinweb}
\textbf{\ta{நாதஸ்வரம்}} (\emph{nātasvaram}) is built from Sanskrit \emph{nāda}, \textquotedblleft{}sound\textquotedblright{} (CDIAL 7052), and \emph{svara}, \textquotedblleft{}a musical note\textquotedblright{} (CDIAL 13908).
\end{cousinweb}

\begin{grammarlens}[title={What you've built}]
One more word for saying what happens if.
\end{grammarlens}

\subsection*{Your turn}
\begin{itemize}
\raggedright
  \item \emph{Say it:} \emph{nātasvaram}
  \item \emph{Say it:} \emph{nātasvaram}, once more
  \item \emph{Say it:} \emph{nātasvaram illāmal tirumaṇam naṭakkātu eṉṟu pāṭṭi colkiṟāḷ.} — Grandmother says a wedding cannot happen without the nagaswaram.
  \item \emph{Recall:} say the Tamil for a guitar, then the Tamil for a violin, then say \emph{nātasvaram} again
\end{itemize}

\begin{tcolorbox}[breakable,colback=teal!4,colframe=teal!35!black,title={Before you move on}]
What does \ta{நாதஸ்வரம்} mean? (\textquotedblleft{}A nagaswaram\textquotedblright{}.) Say it once more.
\end{tcolorbox}
\section[\texorpdfstring{\ta{மேளம்} (mēḷam)}{மேளம் (mēḷam)}]{\ta{மேளம்} (mēḷam) — a drum; a drum band}

\label{lesson:TA-C558-melam}



\begin{quote}
Before the new one: say the Tamil for a violin, then the Tamil for a nagaswaram.
\end{quote}

\subsection*{You'll want to know: \ta{மேளம்}}
\textbf{\ta{மேளம்}} — \emph{mēḷam} — \textquotedblleft{}a drum; a drum band\textquotedblright{}.

\begin{quote}
\ta{மேளம்} \ta{கிடைக்காவிட்டால்} \ta{பதிவு} \ta{செய்த} \ta{இசையைப்} \ta{போடுவோம்}.
\end{quote}

\begin{quote}
\emph{mēḷam kiṭaikkāviṭṭāl pativu ceyta icaiyaip pōṭuvōm.}
\end{quote}

\begin{quote}
If we can't get drummers, we will play recorded music.
\end{quote}

At a wedding, \ta{மேளம்} often means the whole band of drums and pipes.

\begin{cousinweb}
The origin of \textbf{\ta{மேளம்}} (\emph{mēḷam}) is not settled: no dictionary we can check gives its source.
\end{cousinweb}

\begin{grammarlens}[title={What you've built}]
One more word for saying what happens if.
\end{grammarlens}

\subsection*{Your turn}
\begin{itemize}
\raggedright
  \item \emph{Say it:} \emph{mēḷam}
  \item \emph{Say it:} \emph{mēḷam}, once more
  \item \emph{Say it:} \emph{mēḷam kiṭaikkāviṭṭāl pativu ceyta icaiyaip pōṭuvōm.} — If we can't get drummers, we will play recorded music.
  \item \emph{Recall:} say the Tamil for a violin, then the Tamil for a nagaswaram, then say \emph{mēḷam} again
\end{itemize}

\begin{tcolorbox}[breakable,colback=teal!4,colframe=teal!35!black,title={Before you move on}]
What does \ta{மேளம்} mean? (\textquotedblleft{}A drum; a drum band\textquotedblright{}.) Say it once more.
\end{tcolorbox}
\section[\texorpdfstring{\ta{ராகம்} (rākam)}{ராகம் (rākam)}]{\ta{ராகம்} (rākam) — a raga; a melody}

\label{lesson:TA-C558-rakam}



\begin{quote}
Before the new one: say the Tamil for a nagaswaram, then the Tamil for a drum.
\end{quote}

\subsection*{You'll want to know: \ta{ராகம்}}
\textbf{\ta{ராகம்}} — \emph{rākam} — \textquotedblleft{}a raga; a melody\textquotedblright{}.

\begin{quote}
\ta{அவள்} \ta{பாடிய} \ta{ராகம்} \ta{எனக்குத்} \ta{தெரியவில்லை}.
\end{quote}

\begin{quote}
\emph{avaḷ pāṭiya rākam eṉakkut teriyavillai.}
\end{quote}

\begin{quote}
I didn't recognise the raga she sang.
\end{quote}

A melodic framework of Indian classical music. A -\ta{அம்} noun: \ta{ராகத்தில்}.

\begin{cousinweb}
\textbf{\ta{ராகம்}} (\emph{rākam}) is borrowed from Sanskrit \emph{rāga}, \textquotedblleft{}colour, passion\textquotedblright{}, and in music \textquotedblleft{}a melodic mode\textquotedblright{} (Monier-Williams).
\end{cousinweb}

\begin{grammarlens}[title={What you've built}]
One more word for saying what happens if.
\end{grammarlens}

\subsection*{Your turn}
\begin{itemize}
\raggedright
  \item \emph{Say it:} \emph{rākam}
  \item \emph{Say it:} \emph{rākam}, once more
  \item \emph{Say it:} \emph{avaḷ pāṭiya rākam eṉakkut teriyavillai.} — I didn't recognise the raga she sang.
  \item \emph{Recall:} say the Tamil for a nagaswaram, then the Tamil for a drum, then say \emph{rākam} again
\end{itemize}

\begin{tcolorbox}[breakable,colback=teal!4,colframe=teal!35!black,title={Before you move on}]
What does \ta{ராகம்} mean? (\textquotedblleft{}A raga; a melody\textquotedblright{}.) Say it once more.
\end{tcolorbox}
\section[\texorpdfstring{\ta{தாளம்} (tāḷam)}{தாளம் (tāḷam)}]{\ta{தாளம்} (tāḷam) — rhythm; a beat}

\label{lesson:TA-C558-talam}



\begin{quote}
Before the new one: say the Tamil for a drum, then the Tamil for a raga.
\end{quote}

\subsection*{You'll want to know: \ta{தாளம்}}
\textbf{\ta{தாளம்}} — \emph{tāḷam} — \textquotedblleft{}rhythm; a beat\textquotedblright{}.

\begin{quote}
\ta{தாளம்} \ta{தவறினால்} \ta{பாடல்} \ta{கெட்டுவிடும்}.
\end{quote}

\begin{quote}
\emph{tāḷam tavaṟiṉāl pāṭal keṭṭuviṭum.}
\end{quote}

\begin{quote}
If the beat goes wrong, the song is spoiled.
\end{quote}

One keeps time with the hand: \ta{தாளம்} \ta{போடு}.

\begin{cousinweb}
\textbf{\ta{தாளம்}} (\emph{tāḷam}) is borrowed from Sanskrit \emph{tāla}, \textquotedblleft{}a cymbal\textquotedblright{} (CDIAL 5801), and in music \textquotedblleft{}a time measure\textquotedblright{}.
\end{cousinweb}

\begin{grammarlens}[title={What you've built}]
One more word for saying what happens if.
\end{grammarlens}

\subsection*{Your turn}
\begin{itemize}
\raggedright
  \item \emph{Say it:} \emph{tāḷam}
  \item \emph{Say it:} \emph{tāḷam}, once more
  \item \emph{Say it:} \emph{tāḷam tavaṟiṉāl pāṭal keṭṭuviṭum.} — If the beat goes wrong, the song is spoiled.
  \item \emph{Recall:} say the Tamil for a drum, then the Tamil for a raga, then say \emph{tāḷam} again
\end{itemize}

\begin{tcolorbox}[breakable,colback=teal!4,colframe=teal!35!black,title={Before you move on}]
What does \ta{தாளம்} mean? (\textquotedblleft{}Rhythm; a beat\textquotedblright{}.) Say it once more.
\end{tcolorbox}
\section[\texorpdfstring{\ta{பியானோ} (piyāṉō)}{பியானோ (piyāṉō)}]{\ta{பியானோ} (piyāṉō) — a piano}

\label{lesson:TA-C558-piyano}



\begin{quote}
Before the new one: say the Tamil for a raga, then the Tamil for rhythm.
\end{quote}

\subsection*{You'll want to know: \ta{பியானோ}}
\textbf{\ta{பியானோ}} — \emph{piyāṉō} — \textquotedblleft{}a piano\textquotedblright{}.

\begin{quote}
\ta{மண்டபத்தில்} \ta{பியானோ} \ta{இருந்தால்} \ta{அவள்} \ta{வாசிப்பாள்}.
\end{quote}

\begin{quote}
\emph{maṇṭapattil piyāṉō iruntāl avaḷ vācippāḷ.}
\end{quote}

\begin{quote}
If there is a piano in the hall, she will play.
\end{quote}

Played with \ta{வாசி}: \ta{பியானோ} \ta{வாசி}.

\begin{cousinweb}
\textbf{\ta{பியானோ}} (\emph{piyāṉō}) is a loan from English \emph{piano}.
\end{cousinweb}

\begin{grammarlens}[title={What you've built}]
One more word for saying what happens if.
\end{grammarlens}

\subsection*{Your turn}
\begin{itemize}
\raggedright
  \item \emph{Say it:} \emph{piyāṉō}
  \item \emph{Say it:} \emph{piyāṉō}, once more
  \item \emph{Say it:} \emph{maṇṭapattil piyāṉō iruntāl avaḷ vācippāḷ.} — If there is a piano in the hall, she will play.
  \item \emph{Recall:} say the Tamil for a raga, then the Tamil for rhythm, then say \emph{piyāṉō} again
\end{itemize}

\begin{tcolorbox}[breakable,colback=teal!4,colframe=teal!35!black,title={Before you move on}]
What does \ta{பியானோ} mean? (\textquotedblleft{}A piano\textquotedblright{}.) Say it once more.
\end{tcolorbox}
\section[(dialogue)]{Looking back at plans, risks and rules}

\label{lesson:TA-R558-review-552-558}



\begin{quote}
Say the Tamil for a fire extinguisher and a gas cylinder.
\end{quote}

\subsection*{Your turn}
Say these aloud:

\begin{itemize}
\raggedright
  \item a fire extinguisher, a gas cylinder, a buoy, to speed up, time to act
  \item temporary, permanent, essential, an alternative, approximately
  \item at most, as far as possible, trekking, to go fishing, a puppet show
  \item applause, an auditorium, an award, a review, a story character
  \item a short story, literature, a craft, weaving, gold thread
  \item cotton, wool, a dye, a guitar, a violin
  \item a nagaswaram, a drum, a raga, rhythm, a piano
\end{itemize}

\begin{tcolorbox}[breakable,colback=teal!4,colframe=teal!35!black,title={Before you move on}]
A fire extinguisher? (\textbf{\ta{தீயணைப்பான்}}, \emph{tīyaṇaippāṉ}.) A gas cylinder? (\textbf{\ta{சிலிண்டர்}}, \emph{ciliṇṭar}.) A buoy? (\textbf{\ta{மிதவை}}, \emph{mitavai}.) To speed up? (\textbf{\ta{துரிதப்படுத்து}}, \emph{turitappaṭuttu}.) Time to act? (\textbf{\ta{அவகாசம்}}, \emph{avakācam}.) Temporary? (\textbf{\ta{தற்காலிக}}, \emph{taṟkālika}.) Permanent? (\textbf{\ta{நிரந்தர}}, \emph{nirantara}.) Essential? (\textbf{\ta{இன்றியமையாத}}, \emph{iṉṟiyamaiyāta}.) An alternative? (\textbf{\ta{மாற்றுவழி}}, \emph{māṟṟuvaḻi}.) Approximately? (\textbf{\ta{தோராயமாக}}, \emph{tōrāyamāka}.) At most? (\textbf{\ta{அதிகபட்சம்}}, \emph{atikapaṭcam}.) As far as possible? (\textbf{\ta{முடிந்தவரை}}, \emph{muṭintavarai}.) Trekking? (\textbf{\ta{மலையேற்றம்}}, \emph{malaiyēṟṟam}.) To go fishing? (\textbf{\ta{மீன்பிடி}}, \emph{mīṉpiṭi}.) A puppet show? (\textbf{\ta{பொம்மலாட்டம்}}, \emph{pommalāṭṭam}.) Applause? (\textbf{\ta{கைத்தட்டல்}}, \emph{kaittaṭṭal}.) An auditorium? (\textbf{\ta{அரங்கம்}}, \emph{araṅkam}.) An award? (\textbf{\ta{விருது}}, \emph{virutu}.) A review? (\textbf{\ta{விமர்சனம்}}, \emph{vimarcaṉam}.) A story character? (\textbf{\ta{கதாபாத்திரம்}}, \emph{katāpāttiram}.) A short story? (\textbf{\ta{சிறுகதை}}, \emph{ciṟukatai}.) Literature? (\textbf{\ta{இலக்கியம்}}, \emph{ilakkiyam}.) A craft? (\textbf{\ta{கைவினை}}, \emph{kaiviṉai}.) Weaving? (\textbf{\ta{நெசவு}}, \emph{necavu}.) Gold thread? (\textbf{\ta{சரிகை}}, \emph{carikai}.) Cotton? (\textbf{\ta{பருத்தி}}, \emph{parutti}.) Wool? (\textbf{\ta{கம்பளி}}, \emph{kampaḷi}.) A dye? (\textbf{\ta{சாயம்}}, \emph{cāyam}.) A guitar? (\textbf{\ta{கிட்டார்}}, \emph{kiṭṭār}.) A violin? (\textbf{\ta{வயலின்}}, \emph{vayaliṉ}.) A nagaswaram? (\textbf{\ta{நாதஸ்வரம்}}, \emph{nātasvaram}.) A drum? (\textbf{\ta{மேளம்}}, \emph{mēḷam}.) A raga? (\textbf{\ta{ராகம்}}, \emph{rākam}.) Rhythm? (\textbf{\ta{தாளம்}}, \emph{tāḷam}.) A piano? (\textbf{\ta{பியானோ}}, \emph{piyāṉō}.)
\end{tcolorbox}
\section[(dialogue)]{Booking Musicians — once more}

\label{lesson:TA-R558-review-final-558}



\begin{quote}
Say the Tamil for a nagaswaram and a drum.
\end{quote}

\subsection*{Your turn}
\begin{itemize}
\raggedright
  \item \emph{Say it:} a nagaswaram, a drum, a raga, rhythm, a piano
\end{itemize}

\begin{tcolorbox}[breakable,colback=teal!4,colframe=teal!35!black,title={Before you move on}]
A nagaswaram? (\textbf{\ta{நாதஸ்வரம்}}, \emph{nātasvaram}.) A drum? (\textbf{\ta{மேளம்}}, \emph{mēḷam}.) A raga? (\textbf{\ta{ராகம்}}, \emph{rākam}.) Rhythm? (\textbf{\ta{தாளம்}}, \emph{tāḷam}.) A piano? (\textbf{\ta{பியானோ}}, \emph{piyāṉō}.)
\end{tcolorbox}
